\begin{pastexam}[label={px:gw3-106-ds-04}]{關務三等資訊處理 106 年 第四題}
（申論；25 分）

有一雜湊表格（hash table）包含 11 個桶（buckets），位址編號由 0 至 10，每個桶有一個槽（slot）。雜湊函數 h 的定義為
h（key）= key \% 11（註：a\%b 表示 a 除以 b 的餘數）。當有碰撞（collision）發生時，採用線性探測（linear probing）
解決碰撞問題。從空的雜湊表格開始，依序加入 10 個整數 5, 51, 23, 68, 12, 36, 6, 30, 32, 10。
\begin{enumerate}[leftmargin=2.6em, labelsep=0.4em]
\item[㈠] 請繪出加入 10 個整數後的雜湊表格。（15 分）
\item[㈡] 欲在此雜湊表格中尋找資料值 35，請說明須經過幾次的資料值比對，才能確定資料值 35 不在此雜湊表格中。（10 分）
\end{enumerate}
\end{pastexam}

\begin{pastsoln}
答案：㈠ 見下表；㈡ $7$ 次資料值比對。

\textbf{㈠} 線性探測（資結書第~12.3~節）：先放 $h(\mathit{key})$，被占用就往下一格 $(\cdot + 1) \bmod 11$，直到空格。
\begin{center}
\small
\begin{tabular}{ccl}
\toprule
加入 & $h = \mathit{key} \bmod 11$ & 探測過程 \\
\midrule
5 & 5 & 位址 5 空，放入 \\
51 & 7 & 位址 7 空，放入 \\
23 & 1 & 位址 1 空，放入 \\
68 & 2 & 位址 2 空，放入 \\
12 & 1 & 1（23）、2（68）占用，放入 \textbf{3} \\
36 & 3 & 3（12）占用，放入 \textbf{4} \\
6 & 6 & 位址 6 空，放入 \\
30 & 8 & 位址 8 空，放入 \\
32 & 10 & 位址 10 空，放入 \\
10 & 10 & 10（32）占用，繞回 \textbf{0}，放入 \\
\bottomrule
\end{tabular}
\end{center}
（$51 = 4 \cdot 11 + 7$，$68 = 6 \cdot 11 + 2$，$36 = 3 \cdot 11 + 3$，$30 = 2 \cdot 11 + 8$，$32 = 2 \cdot 11 + 10$。）
\begin{center}
\small
\begin{tabular}{c|ccccccccccc}
位址 & 0 & 1 & 2 & 3 & 4 & 5 & 6 & 7 & 8 & 9 & 10 \\
\hline
資料 & 10 & 23 & 68 & 12 & 36 & 5 & 6 & 51 & 30 & & 32 \\
\end{tabular}
\end{center}
只有位址 $9$ 是空的。

\textbf{㈡} 搜尋走同樣的探測順序：$h(35) = 35 \bmod 11 = 2$，從位址 $2$ 開始往下比，遇到空格才能確定不存在（空格
之後不可能有 35，因為當初插入 35 的話會停在這個空格或更前面）。
\begin{center}
\small
\begin{tabular}{cccc}
\toprule
第幾次比對 & 位址 & 資料 & 結果 \\
\midrule
1 & 2 & 68 & $\ne 35$ \\
2 & 3 & 12 & $\ne 35$ \\
3 & 4 & 36 & $\ne 35$ \\
4 & 5 & 5 & $\ne 35$ \\
5 & 6 & 6 & $\ne 35$ \\
6 & 7 & 51 & $\ne 35$ \\
7 & 8 & 30 & $\ne 35$ \\
--- & 9 & 空 & 停止，35 不在表中 \\
\bottomrule
\end{tabular}
\end{center}
須和 $7$ 個資料值比對；連同最後檢查到空的位址 $9$，共探測 $8$ 個位址。 線性探測的群聚現象在這裡很明顯：位址 $10$ 和
$0$--$8$ 繞成一長段，任何落在這段裡的搜尋失敗都要走到位址 $9$ 才停。 （若表中沒有空格，就要繞一圈比完 $11$ 格才能
確定不存在；這也是開放定址的表不能放滿的原因。）
\end{pastsoln}
